% ============================================================
% SCHEDE PG EXTRA — CASA DORIA
% Due personaggi aggiuntivi per tavolo alternativo / quinta sessione
%
% @rpg-schema: <https://rpg-schema.org/instances/genova1507/crew-doria>
%   fitd:hasOptionalMember
%     <.../pg-elena-doria> ,
%     <.../pg-luca-grimbert> .
% ============================================================

\chapter{Schede PG Extra --- Casa Doria}
\index{schede PG!Doria!extra}

% ════════════════════════════════════════════════════════════
% PG EXTRA 1 — ELENA DORIA
% @rpg-schema: <https://rpg-schema.org/instances/genova1507/pg-elena-doria>
%   a fitd:Character ;
%   rdfs:label "Elena Doria" ;
%   fitd:crew <.../crew-doria> ;
%   fitd:role "La Maestra Calafata / L'Arsenale" ;
%   fitd:action [ fitd:actionName "Tinker" ; fitd:dots 4 ] ;
%   fitd:action [ fitd:actionName "Wreck" ; fitd:dots 3 ] ;
%   fitd:action [ fitd:actionName "Command" ; fitd:dots 2 ] .
% ════════════════════════════════════════════════════════════

\pgheader[DoriaBlue]{Elena Doria}{La Maestra Calafata / L'Arsenale}{Doria}
\pgquote{Una nave che non galleggia non è una nave. È un problema.}

\section*{Background \& Motivazione}

\begin{rulebox}{}
  \textbf{Background:} Quarantaquattro anni, figlia illegittima di un cugino
  minore della famiglia. Ha imparato il mestiere di calafato nell'arsenale di Prè
  a fianco del padre prima ancora di sapere leggere. Ora è la caposquadra
  de facto dell'intero cantiere: nessuna delle quattro galee Doria entra in acqua
  senza il suo sigillo. I marinai la chiamano semplicemente «la Maestra».
  La famiglia ufficialmente la ignora. Ufficiosa­mente la consulta su tutto.

  \medskip\noindent\textcolor{DoriaBlue}{\textbf{Motivazione:}} Elena non ha
  ideologie. Ha galee. Se i francesi vincono, le galee Doria vengono confiscate
  o affondate — e con esse trent'anni di lavoro suo e dei suoi operai. Questo
  non è accettabile. La rivolta è quindi, per lei, una questione professionale.
\end{rulebox}

\section*{Attributi \& Azioni}

\begin{multicols}{2}

\begin{attributoblock}[DoriaBlue]{INSIGHT --- Mente}
  \azione{Caccia}{Hunt}{1}{Individua difetti strutturali, materiali di contrabbando}
  \azione{Studio}{Study}{2}{Architettura navale, contratti d'appalto, materiali}
  \azione{Osserva}{Survey}{3}{Vede immediatamente cosa non va in una nave o in una persona}
  \azione{Ingegno}{Tinker}{4}{Costruzione, sabotaggio, riparazione di qualsiasi imbarcazione}
\end{attributoblock}

\columnbreak

\begin{attributoblock}[DoriaBlue]{PROWESS --- Corpo}
  \azione{Finezza}{Finesse}{2}{Mani precise da artigiana — sa usare un coltello}
  \azione{Ombra}{Prowl}{1}{Conosce ogni angolo dell'arsenale al buio}
  \azione{Combattimento}{Skirmish}{2}{Ruvida, efficace, senza stile}
  \azione{Forza bruta}{Wreck}{3}{Trent'anni di martello e scalpello}
\end{attributoblock}

\end{multicols}

\begin{attributoblock}[DoriaBlue]{RESOLVE --- Spirito}
  \begin{multicols}{2}
    \azione{Percezione}{Attune}{1}{Sente quando qualcosa nell'arsenale non torna}
    \azione{Comando}{Command}{2}{Comanda gli operai con autorità guadagnata, non ereditata}
    \azione{Relazioni}{Consort}{2}{Maestranze, fornitori, capitani di porto}
    \azione{Persuasione}{Sway}{1}{Diretta al punto dell'impazienza}
  \end{multicols}
\end{attributoblock}

\medskip
\begin{pgclockbox}{Stress \hfill \clock{9}{0}}
  \small\textit{Traumi: $\square$ Freddo \quad $\square$ Duro \quad $\square$ Ossessionato \quad $\square$ Paranoico \quad $\square$ Irrazionale \quad $\square$ Spezzato}
\end{pgclockbox}

\section*{Abilità Speciale}

% @rpg-schema: <.../pg-elena-doria>
%   fitd:specialAbility [ rdfs:label "Sapere dell'Arsenale" ] .

\begin{doriabox}{Sapere dell'Arsenale}
  Elena conosce ogni nave nel porto di Genova — scafo, carico, stato di
  manutenzione, equipaggio reale. Può dichiarare una verità tecnica su
  qualsiasi imbarcazione (è lenta, ha una falla, può essere affondata con
  una modifica specifica) e quella verità è sempre accurata.

  \medskip\noindent\textit{Nessun tiro. Una volta per scena. Se la verità
  dichiarata richiede azione, l'azione richiede un tiro appropriato.}
\end{doriabox}

\section*{Legami}

\begin{table}[h]
\centering\renewcommand{\arraystretch}{1.4}
\begin{tabularx}{\linewidth}{@{}L{2.4cm} X@{}}
  \toprule\rowcolor{DoriaBlue}
  \textcolor{white}{\textbf{PG}} & \textcolor{white}{\textbf{Legame}} \\
  \midrule
  Ottaviano & Mi rispetta. Non mi ha mai chiamato figlia, ma mi ha insegnato
             a nuotare da bambina. È sufficiente. \\
  \rowcolor{DoriaBlue!8}
  Ser Fulvio & Sa che posso sabotare una galea senza che nessuno se ne accorga
              prima di tre giorni di navigazione. Questo ci rende utili a vicenda. \\
  Bianca & L'unica in famiglia che mi ha chiesto il mio nome per intero.
           La rispetto per questo. \\
  \rowcolor{DoriaBlue!8}
  Giacomo & Mi fornisce i materiali che non posso ordinare ufficialmente.
            Non chiedo da dove vengono. \\
  \bottomrule
\end{tabularx}
\end{table}

\section*{Equipaggiamento}
\begin{goldlist}
  \item Martello da calafato con manico in corno di bue --- strumento e arma
  \item Quaderno tecnico con i difetti strutturali delle galee francesi nel porto
  \item Chiave dell'arsenale e di tre depositi di materiali
  \item Sacchetto di stoppa impregnata di pece --- appiccica il fuoco in trenta secondi
\end{goldlist}

\clearpage

\begin{secretbox}
  \textbf{Segreto:} Elena sa che due delle galee Doria sono ipotecate --- non
  perché Ottaviano glielo abbia detto, ma perché ha visto i documenti passare
  per l'arsenale. Sa anche chi è il banchiere creditore. E sa che quel banchiere
  ha un magazzino nel porto con materiali rubati dall'arsenale stesso negli
  ultimi tre anni. Ha prove. Non le ha ancora usate.

  \medskip\noindent\textcolor{SecretRed}{\textbf{Agenda:}}
  \textit{Garantire che le galee rimangano in mani genovesi a qualsiasi costo.
  Se la situazione lo richiede, è disposta a ricattare il banchiere creditore
  con le prove che ha — anche senza l'autorizzazione di Ottaviano.}
\end{secretbox}

\clearpage

% ════════════════════════════════════════════════════════════
% PG EXTRA 2 — LUCA GRIMBERT
% @rpg-schema: <https://rpg-schema.org/instances/genova1507/pg-luca-grimbert>
%   a fitd:Character ;
%   rdfs:label "Luca Grimbert" ;
%   fitd:crew <.../crew-doria> ;
%   fitd:role "L'Ufficiale Francese / Il Traditore" ;
%   fitd:action [ fitd:actionName "Survey" ; fitd:dots 3 ] ;
%   fitd:action [ fitd:actionName "Consort" ; fitd:dots 3 ] ;
%   fitd:action [ fitd:actionName "Sway" ; fitd:dots 3 ] .
% ════════════════════════════════════════════════════════════

\pgheader[DoriaBlue]{Luca Grimbert}{L'Ufficiale Francese / Il Traditore}{Doria}
\pgquote{Servo il re di Francia. Ma il re di Francia non sa dove sono stanotte.}

\section*{Background \& Motivazione}

\begin{rulebox}{}
  \textbf{Background:} Trentuno anni, figlio di un mercante provenzale e di
  una genovese. Ufficiale subalterno della guarnigione francese al Castelletto,
  incaricato della logistica dei rifornimenti. Da diciotto mesi passa
  informazioni ai Doria — non per ideologia, ma perché Andrea Doria gli ha
  salvato il fratello da un processo per contrabbando e il debito non è stato
  mai formalizzato né mai dimenticato.

  \medskip\noindent\textcolor{DoriaBlue}{\textbf{Motivazione:}} Sopravvivere
  alla notte. Luca sa che la rivolta è imminente e sa che essere un ufficiale
  francese a Genova al mattino dopo sarà pericoloso. Ha bisogno che i Doria
  vincano abbastanza da proteggerlo — ma non abbastanza da non aver più bisogno
  di lui.
\end{rulebox}

\section*{Attributi \& Azioni}

\begin{multicols}{2}

\begin{attributoblock}[DoriaBlue]{INSIGHT --- Mente}
  \azione{Caccia}{Hunt}{2}{Logistica militare, rotazioni di guardia, movimenti truppe}
  \azione{Studio}{Study}{2}{Francese, italiano, latino, protocollo militare}
  \azione{Osserva}{Survey}{3}{Addestrato a rilevare anomalie — le sue e quelle altrui}
  \azione{Ingegno}{Tinker}{1}{Armi, fortificationi leggere}
\end{attributoblock}

\columnbreak

\begin{attributoblock}[DoriaBlue]{PROWESS --- Corpo}
  \azione{Finezza}{Finesse}{2}{Spada da ufficiale, movimenti calibrati}
  \azione{Ombra}{Prowl}{3}{Muoversi tra le due fazioni senza essere visto}
  \azione{Combattimento}{Skirmish}{2}{Addestramento militare standard}
  \azione{Forza bruta}{Wreck}{1}{Non la sua specialità}
\end{attributoblock}

\end{multicols}

\begin{attributoblock}[DoriaBlue]{RESOLVE --- Spirito}
  \begin{multicols}{2}
    \azione{Percezione}{Attune}{2}{Sente quando una copertura sta cedendo}
    \azione{Comando}{Command}{1}{Comanda i subalterni francesi con disagio crescente}
    \azione{Relazioni}{Consort}{3}{Frequenta entrambi i mondi — ufficiali francesi e nobili genovesi}
    \azione{Persuasione}{Sway}{3}{Convince con la plausibilità, non con la forza}
  \end{multicols}
\end{attributoblock}

\medskip
\begin{pgclockbox}{Stress \hfill \clock{9}{0}}
  \small\textit{Traumi: $\square$ Freddo \quad $\square$ Duro \quad $\square$ Ossessionato \quad $\square$ Paranoico \quad $\square$ Irrazionale \quad $\square$ Spezzato}
\end{pgclockbox}

\section*{Abilità Speciale}

% @rpg-schema: <.../pg-luca-grimbert>
%   fitd:specialAbility [ rdfs:label "Ordini Contraffatti" ] .

\begin{doriabox}{Ordini Contraffatti}
  Luca può produrre ordini militari francesi plausibili — ritiro di pattuglie,
  riallocazione di guardie, autorizzazione al transito — che funzionano
  su qualsiasi soldato francese di rango inferiore al suo.
  Gli ordini reggono finché non vengono verificati con la catena di comando.

  \medskip\noindent\textit{Tiro su Studio per ordini complessi o per truppe
  sospettose. Ordini semplici non richiedono tiro. Ogni uso aumenta il rischio
  che la sua copertura venga scoperta — il GM traccia la pressione.}
\end{doriabox}

\section*{Legami}

\begin{table}[h]
\centering\renewcommand{\arraystretch}{1.4}
\begin{tabularx}{\linewidth}{@{}L{2.4cm} X@{}}
  \toprule\rowcolor{DoriaBlue}
  \textcolor{white}{\textbf{PG}} & \textcolor{white}{\textbf{Legame}} \\
  \midrule
  Ottaviano & L'Ammiraglio mi tratta come uno strumento. Preferisco così.
             Gli strumenti non vengono sacrificati per principio. \\
  \rowcolor{DoriaBlue!8}
  Ser Fulvio & Mi gestisce direttamente. Sa più di me di quante persone
              mi conoscano in entrambi i campi. Questo mi spaventa. \\
  Bianca & Ha letto i miei rapporti prima ancora che li consegnassi.
           Non so come. Non ho chiesto. \\
  \rowcolor{DoriaBlue!8}
  Elena & L'unica che non mi chiede mai informazioni. Mi guarda come
         se sapesse già tutto. Forse è così. \\
  \bottomrule
\end{tabularx}
\end{table}

\section*{Equipaggiamento}
\begin{goldlist}
  \item Uniforme da ufficiale francese completa --- include sigillo e documenti di identità
  \item Abiti borghesi genovesi per muoversi fuori dalla guarnigione
  \item Chiave del magazzino dei rifornimenti al Castelletto
  \item Lettera del fratello --- tiene in tasca. Non sa perché.
\end{goldlist}

\clearpage

\begin{secretbox}
  \textbf{Segreto:} Il comandante della guarnigione ha ricevuto ieri sera una
  soffiata: c'è un traditore tra gli ufficiali. Non sa chi. Luca lo sa perché
  era presente alla riunione. Ha quarantotto ore, forse meno, prima che le
  indagini lo raggiungano. Non lo ha detto ai Doria perché non vuole che
  accelerino i piani prima che lui abbia una via di uscita garantita.

  \medskip\noindent\textcolor{SecretRed}{\textbf{Agenda:}}
  \textit{Ottenere una garanzia scritta di protezione dalla famiglia Doria
  prima della fine della sessione. Finché non ce l'ha, passa solo le
  informazioni strettamente necessarie a mantenersi utile — non quelle
  che li farebbero vincere troppo in fretta.}
\end{secretbox}
